\section{资源管理、增长与失效}

\begin{frame}{\texttt{vector}替我们管理什么？}
\begin{center}
\begin{tikzpicture}[>=Stealth,font=\small,
box/.style={draw=VSPPrimary,rounded corners=2pt,minimum height=1.0cm,align=center}]
\node[box,fill=VSPPrimary!7,minimum width=3.2cm] (obj) {\texttt{std::vector}对象\\记录大小、容量等资源状态};
\node[box,fill=VSPPrimary!12,minimum width=5.2cm,right=2.0cm of obj] (buf) {连续元素存储\\增长时可能搬迁，析构时释放};
\draw[->,very thick,VSPPrimary] (obj) -- node[above]{拥有并管理} (buf);
\end{tikzpicture}
\end{center}
\begin{itemize}
\item 局部\texttt{vector}对象离开作用域时自动析构并释放元素存储；这种对象生命周期绑定资源的思路称为\textbf{RAII}。
\item 不需要也不应该对\texttt{vector}的元素缓冲区调用\texttt{free}。
\end{itemize}
\alert{局部\texttt{vector}对象本身与它管理的元素存储是两处不同的内存：对象随作用域自动销毁，元素存储由对象在析构时释放。}
\end{frame}

\begin{frame}{\texttt{size}与\texttt{capacity}不是一回事}
\begin{center}
\begin{tikzpicture}[font=\scriptsize]
\foreach \i in {0,...,7} {
  \pgfmathsetmacro{\x}{\i*1.05}
  \ifnum\i<4
    \filldraw[fill=VSPPrimary!14,draw=VSPPrimary] (\x,0) rectangle (\x+1.05,0.85);
  \else
    \filldraw[fill=black!2,draw=black!35] (\x,0) rectangle (\x+1.05,0.85);
  \fi
}
\node at (2.1,-0.45) {\texttt{size() == 4}：已有元素};
\node at (6.3,-0.45) {\texttt{capacity() == 8}：当前存储可容纳元素};
\draw[<->,VSPPrimary,thick] (0,1.15) -- node[above]{已有元素} (4.2,1.15);
\draw[<->,black!55,thick] (0,1.55) -- node[above]{当前分配容量} (8.4,1.55);
\end{tikzpicture}
\end{center}
\begin{itemize}
\item \texttt{size()}是元素个数；\texttt{capacity()}是无需重新分配即可容纳的元素个数。
\item 恒有\texttt{size() <= capacity()}；容量按元素计，不是字节数。
\item 容量大于大小不代表额外元素已经存在，不能通过下标访问它们。
\end{itemize}
\end{frame}

\begin{frame}{为什么几何增长能做到均摊常数？}
\begin{columns}[T]
\begin{column}{0.47\textwidth}
\textbf{每次只增加1个容量}
\medskip

第$i+1$个元素到来时，复制原有$i$个元素。连续$n$次追加的总复制量：
\[
\sum_{i=0}^{n-1} i = O(n^2)
\]
\end{column}
\begin{column}{0.49\textwidth}
\textbf{容量按倍数增长}
\medskip

例如容量按$0\to1\to2\to4\to8\to\cdots$增长。扩容时总复制量：
\[
1+2+4+\cdots+2^{m-1}=O(n)
\]
均摊到$n$次\texttt{push\_back}，每次为$O(1)$。
\end{column}
\end{columns}
\medskip
\begin{itemize}
\item 单次扩容仍然昂贵；“均摊$O(1)$”描述一串操作的总成本。
\item 具体增长因子由实现选择，标准并不规定必须为2倍。
\end{itemize}
\end{frame}

\begin{frame}{观察实际增长：输出随实现而变化}
\begin{Code}*{cpp}\lstinputlisting{code/growth.cpp}\end{Code}
\begin{itemize}
\item 每次追加后打印\texttt{size}与\texttt{capacity}，观察何时发生容量增长。
\item 可检查不变量：大小依次为$1$到$17$，容量不减且始终不小于大小。
\end{itemize}
\alert{不要把某次运行得到的容量序列写成所有编译器都保证的规则。}
\end{frame}

\begin{frame}{已知大致元素个数：\texttt{reserve()}}
\begin{Code}*{cpp}\lstinputlisting{code/vector-reserve.inc}\end{Code}
\begin{itemize}
\item \texttt{reserve(n)}请求至少容纳$n$个元素的容量，可以减少后续重新分配。
\item 它\textbf{不改变大小}；示例中\texttt{v}仍为空，\texttt{v[0]}与\texttt{v.back()}仍是未定义行为。
\item 需要创建$n$个值初始化元素时使用\texttt{resize(n)}或相应构造函数，而不是\texttt{reserve(n)}。
\end{itemize}
\end{frame}

\begin{frame}{重新分配会使哪些东西失效？}
\begin{Code}*{cpp}\lstinputlisting{code/vector-reference-invalid.inc}\end{Code}
\begin{itemize}
\item 若\texttt{push\_back}触发重新分配，元素被转移到新存储，旧存储被释放。
\item 指向旧元素的\textbf{指针、引用和迭代器全部失效}；再使用就是未定义行为。
\item 是否失效取决于是否重新分配，因此错误可能时隐时现，不能靠一次运行证明安全。
\end{itemize}
\end{frame}

\begin{frame}{遍历时追加：从错误写法到安全写法}
\begin{columns}[T]
\begin{column}{0.47\textwidth}
\textbf{错误：范围for依赖迭代器}
\begin{Code}*{cpp}\lstinputlisting{code/vector-invalidate.inc}\end{Code}
重新分配会使遍历所用的迭代器失效。
\end{column}
\begin{column}{0.49\textwidth}
\textbf{一种安全写法：固定原范围}
\begin{Code}*{cpp}\lstinputlisting{code/vector-safe-update.inc}\end{Code}
先复制当前元素再追加；不保存可能失效的引用。
\end{column}
\end{columns}
\medskip
\alert{不要在基于范围的\texttt{for}遍历容器时改变该容器的大小。}
\end{frame}

\begin{frame}{迁移到C++后，仍然要保留哪些C语言直觉？}
\begin{enumerate}
\item \textbf{类型：}对象知道类型，不代表转换和初始化规则可以忽略。
\item \textbf{边界：}\texttt{size()}给出有效范围；\texttt{[]}不检查，空容器没有首尾元素。
\item \textbf{生命周期：}对象自动析构；保存到对象失效后的指针或引用仍然危险。
\item \textbf{失效：}容器增长可能搬迁元素；不要继续持有旧位置的访问句柄。
\item \textbf{复杂度：}\texttt{push\_back}是均摊$O(1)$，单次扩容仍需移动元素。
\end{enumerate}
\medskip
\alert{优先让标准库拥有资源，把精力留给真正需要推理的边界与算法。}
\end{frame}
